%=============================================================================
% W4i19: SURVEY PREPARATION UPDATE --- c_s = c CONSEQUENCES FOR
% EUCLID / DESI / RUBIN PREDICTIONS
% Wave 4, Iteration 19 | 20 March 2026 | Model: Opus 4.6 | Agent 16
%=============================================================================

\documentclass[11pt,a4paper]{article}
\usepackage[margin=2.2cm]{geometry}
\usepackage{amsmath,amssymb}
\usepackage{booktabs}
\usepackage{enumitem}
\usepackage{float}
\usepackage{xcolor}
\usepackage[most]{tcolorbox}
\usepackage{hyperref}
\usepackage{longtable}
\usepackage{array}

\newcommand{\LCDM}{$\Lambda$CDM}
\newcommand{\DFD}{DFD}
\newcommand{\ee}[1]{\times 10^{#1}}
\newcommand{\kmu}{k_\mu}
\newcommand{\Geff}{G_{\rm eff}}
\newcommand{\Dpsi}{\Delta\psi}
\newcommand{\cs}{c_s}

\newtcolorbox{findbox}[1][]{colback=yellow!5!white, colframe=red!70!black,
  fonttitle=\bfseries, title={#1}}
\newtcolorbox{granitebox}[1][]{colback=green!5,colframe=green!60!black,
  fonttitle=\bfseries,title={#1}}
\newtcolorbox{distbox}[1][]{colback=blue!3,colframe=blue!50!black,
  fonttitle=\bfseries,title={#1}}
\newtcolorbox{crisisbox}[1][]{colback=red!5,colframe=red!70!black,
  fonttitle=\bfseries,title={#1}}
\newtcolorbox{pipelinebox}[1][]{colback=orange!3,colframe=orange!60!black,
  fonttitle=\bfseries,title={#1}}

\title{\textbf{Wave 4, Iteration 19:\\
Survey Preparation --- $c_s = c$ Consequences for\\
Euclid DR1 / DESI Year~3 / Rubin Year~1 Predictions}\\[0.5em]
\large $\psi$ Does Not Cluster $\cdot$ Baryon-Only $\Geff$ Enhancement $\cdot$
Specific Testable Numbers\\[0.3em]
\large Timeline Updates $\cdot$ Euclid DR1 Step-by-Step $\cdot$ Rubin SNe~Ia Statistics}

\author{DFD Research Programme --- Agent 16 (Survey Preparation, Opus 4.6)}
\date{20 March 2026}

\begin{document}
\maketitle

%=============================================================================
\section*{Executive Summary: Top 5 Findings}
%=============================================================================

\begin{findbox}[TOP 5 FINDINGS --- WAVE 4, ITERATION 19]

\begin{enumerate}

\item \textbf{CRITICAL --- $c_s = c$ CHANGES COSMOLOGICAL PREDICTIONS
  BUT PRESERVES THE EUCLID KILLER TEST:}
  With $c_s = c$ confirmed (i17--i18, three independent proofs), the
  $\psi$-field does not cluster on sub-horizon scales. The scalar
  perturbation $\delta\psi$ free-streams at the speed of light, contributing
  zero to the matter power spectrum at $k > H/c$. Structure formation
  proceeds through baryon-only $\Geff$ enhancement, \emph{not} through
  $\psi$-clustering.

  \textbf{Derivation chain}: Action Eq.~(2.11) of v3.2 $\to$ temporal
  completion $K(\Delta)$ $\to$ perturbation equation yields $c_s^2 = c^2$
  (Appendix temporal-completion, three proofs: characteristic analysis,
  WKB, Hamiltonian) $\to$ Jeans length $\lambda_J = c/H \sim$ Hubble
  radius $\to$ $\psi$ perturbations do not grow on sub-Hubble scales
  $\to$ structure formation driven by baryon response to modified
  gravitational potential only.

  \textbf{What survives}: The $\Geff(k) = G_N k^2/(k^2 + \kmu^2)$
  Yukawa form remains correct for the \emph{baryonic growth equation},
  because the quasi-static spatial field equation (Eq.~2.7 of v3.2)
  still governs the gravitational potential sourced by matter.
  The lensing convergence $\kappa$ depends on $\nabla^2\psi_{\rm grav}$,
  which is determined by the quasi-static field equation with its
  $\mu(x)$ nonlinearity.

  \textbf{What changes}: (i)~There is no ``dark matter-like'' $\psi$
  clustering. (ii)~The effective $\Geff(k)$ now describes
  how baryons respond to the modified Poisson equation, not how
  $\psi$ itself clusters. (iii)~The galaxy power spectrum prediction
  changes because only baryonic matter contributes to density fluctuations.
  (iv)~CMB predictions require the full DFD Boltzmann code (mochi\_class)
  because $\psi$ perturbations propagate at $c$ and affect the ISW
  effect differently from a cold dark matter component.

  \textbf{Impact: CRITICAL. The tomographic $S_8$ gradient test SURVIVES
  but its amplitude requires mochi\_class for precise calibration.
  The weekend-viable analysis pipeline from i17 remains valid as a
  discovery tool, but the predicted $\Delta S_8$ amplitude now has
  larger theoretical uncertainty: $\Delta S_8 = 0.05$--$0.10$ (was
  $0.07 \pm 0.02$).}

\item \textbf{CRITICAL --- DESI DR2 (Year~3) IS NOW PUBLIC: 14 Million
  Objects, $f\sigma_8$ Precision 4.7\%, Dynamical DE at 3.1$\sigma$:}
  DESI released DR2 results (March 2025) from three years of operations,
  with BAO measurements from $>$14 million galaxies and quasars spanning
  $0.1 < z < 4.2$. Key results:
  \begin{itemize}[nosep]
  \item BAO distance measurements in 7 redshift bins with sub-percent
    precision.
  \item Combined RSD growth-rate precision: 4.7\% from full-shape
    analysis (4.7M galaxy+quasar spectra, $0.1 < z < 2.1$).
  \item DR1 peculiar velocity consensus: $f\sigma_8 = 0.449 \pm 0.008$.
  \item Dynamical dark energy preferred over \LCDM\ at 3.1$\sigma$
    (DESI BAO + CMB); 2.8--4.2$\sigma$ including SNe.
  \end{itemize}
  \textbf{DFD implications}: DFD predicts BAO distances match \LCDM\
  to $<$0.1\% (distance-bias framework, i14). The DESI dynamical-DE
  preference is \emph{not} a DFD signal --- DFD does not modify the
  expansion history. However, the $f\sigma_8$ measurements are directly
  testable (Section~\ref{sec:desi}).
  \textbf{Impact: CRITICAL. DESI $f\sigma_8$ data enables an immediate
  DFD test with existing public data.}

\item \textbf{NEW --- RUBIN FIRST ALERTS LIVE (Feb 2026), $\sim$250,000
  SNe~Ia/year Expected, 3$\sigma$ Anisotropy Test Feasible with
  Year~1 Data:}
  Rubin Observatory began issuing real-time alerts on 24 February 2026,
  with 800,000 alerts on the first night. Full survey operations yield
  $\sim$250,000 Type~Ia supernovae per year to $z \sim 1.2$.
  Data Preview~2 (science-grade commissioning data) expected July--September
  2026. DR1 expected $\sim$2 years after survey start.

  \textbf{DFD prediction for SNe~Ia psi-screen anisotropy}:
  $C_\ell \propto \ell^{-2}$, peaks at $\ell \sim 3$--15, RMS
  $0.8$--$1.2\%$ in distance modulus residuals. With $c_s = c$ (no
  $\psi$ clustering), the anisotropy arises from the inhomogeneous
  gravitational $\psi$-screen along different lines of sight, sourced
  by the large-scale matter distribution (not by $\psi$ self-clustering).
  The signal is robust to $c_s = c$.

  \textbf{SNe needed for detection} (Section~\ref{sec:rubin}):
  $3\sigma$: $\sim$3,000 SNe~Ia with well-measured distances ($z < 0.5$,
  spectroscopic confirmation). $5\sigma$: $\sim$8,000 SNe~Ia.
  Rubin Year~1 will yield $\sim$10,000--50,000 spectroscopically typed
  SNe~Ia (depending on follow-up resources), making $5\sigma$ detection
  feasible within the first year.
  \textbf{Impact: HIGH. Rubin Year~1 can independently test DFD's
  psi-screen anisotropy at $\geq 5\sigma$.}

\item \textbf{NEW --- EUCLID DR1 CONFIRMED 21 OCTOBER 2026, $\sim$2100
  deg$^2$ PROCESSED (1900 deg$^2$ EXPECTED FINAL):}
  ESA confirms processing of $\sim$2100\,deg$^2$ for DR1, with
  $\sim$1900\,deg$^2$ expected as final coverage ($\sim$500\,deg$^2$
  North, $\sim$1400\,deg$^2$ South). Timeline unchanged from i17.
  Q1 data (63\,deg$^2$, released March 2025) is available for pipeline
  testing now.

  \textbf{$c_s = c$ impact on Euclid analysis}: The weekend-viable
  pipeline (i17, Section~\ref{sec:weekend}) is \emph{unchanged} in
  procedure. The theoretical prediction for $\Delta S_8$ now carries
  larger uncertainty (see Finding~1). The step-by-step plan remains
  valid: the observable is $S_8(\ell_{\rm max})$ as a function of
  maximum multipole, with DFD predicting a rising curve and \LCDM\
  predicting a flat line. The \emph{qualitative} signature (monotonic
  rise) is robust to $c_s = c$; only the \emph{amplitude} requires
  mochi\_class calibration.
  \textbf{Impact: HIGH. No change to analysis strategy; theoretical
  error bars widen but discovery potential unchanged.}

\item \textbf{NEW --- STRATEGIC WINDOW CONFIRMED: NO TEAM TESTING
  YUKAWA $\Geff(k)$ ON CURRENT DATA:}
  Cross-checking the competitive landscape (i17 Section~6) against
  DESI DR2 publications and Euclid Q1 papers: no group has tested
  a Yukawa $\Geff(k) = G_N k^2/(k^2 + \kmu^2)$ profile against
  any dataset. The Euclid Consortium's modified-gravity pipeline
  uses scale-independent $(\mu_{\rm mg}, \Sigma_{\rm mg})$. DESI
  analyses test $w_0 w_a$CDM and standard extensions. The strategic
  window for a first-mover DFD paper remains wide open.
  \textbf{Impact: HIGH. The opportunity cost of delay increases as
  Euclid DR1 approaches.}

\end{enumerate}
\end{findbox}


%=============================================================================
\section{The $c_s = c$ Paradigm Shift: What Changes for Surveys}
\label{sec:cs_impact}
%=============================================================================

\begin{crisisbox}[$c_s = c$ --- WHAT EXACTLY CHANGES]

The dust-branch crisis (i17--i18) established that $c_s = c$ at all
epochs. Here we trace the consequences for each survey observable.

\subsection{Derivation Chain for $c_s = c$}

\begin{enumerate}[nosep]
\item \textbf{Starting point}: Full dynamic action Eq.~(2.11) of v3.2:
\[
S_\psi = \int dt\,d^3x\;\left\{
  \frac{a_*^2}{8\pi G}\left[
    W\!\left(\frac{|\nabla\psi|^2}{a_*^2}\right)
    + K\!\left(\frac{c}{a_0}|\dot\psi{-}\dot\psi_0|\right)
  \right]
  - \frac{c^2}{2}\,\psi(\rho - \bar\rho)
\right\}.
\]

\item \textbf{Temporal kinetic term}: $K(\Delta)$ with $\Delta = (c/a_0)|\dot\psi - \dot\psi_0|$.
  The temporal $\mu$-function $K'(\Delta) = \mu(\Delta)$ has the same
  functional form as the spatial $\mu(x) = x/(1+x)$.

\item \textbf{Perturbation analysis}: Linearize $\psi = \bar\psi + \delta\psi$
  around the homogeneous background. The perturbation equation is:
\[
\ddot{\delta\psi} - c^2 \nabla^2 \delta\psi + m_{\rm eff}^2 \delta\psi
  = -\frac{4\pi G c^2}{a_*^2}\,\mu'(\bar\Delta)\,\delta\rho\,,
\]
  where the propagation speed $c_s^2 = c^2$ arises from the ratio
  of spatial to temporal kinetic coefficients. Three independent proofs
  (characteristic analysis, WKB dispersion, Hamiltonian constraint)
  all give $c_s = c$.

\item \textbf{Consequence}: The Jeans wavenumber
  $k_J = \sqrt{4\pi G \bar\rho}/c_s \sim H/c$ is at the Hubble scale.
  All sub-Hubble perturbations in $\delta\psi$ oscillate and decay
  (free-streaming). $\psi$ does not cluster.
\end{enumerate}

\subsection{What Survives Unchanged}

\begin{itemize}[nosep]
\item \textbf{Quasi-static field equation}: For $k \gg H/c$ (all
  scales relevant to surveys), the spatial field equation
  $\nabla\cdot[\mu(|\nabla\psi|/a_*)\nabla\psi] = -(8\pi G/c^2)\rho$
  is valid. This determines the gravitational potential.
\item \textbf{Weak lensing}: Lensing convergence $\kappa$ depends on
  $\nabla^2(\psi_{\rm grav})$ integrated along the line of sight.
  The quasi-static equation governs $\psi_{\rm grav}$.
  \textbf{Lensing predictions are ROBUST to $c_s = c$.}
\item \textbf{$\Geff(k)$ Yukawa form}: Linearizing the quasi-static
  field equation around the homogeneous background gives
  $\Geff(k)/G_N = k^2/(k^2 + \kmu^2)$ with $\kmu = 0.01$\,Mpc$^{-1}$.
  This is a property of the spatial sector, not the temporal sector.
  \textbf{Unchanged.}
\item \textbf{Tomographic $S_8$ gradient}: The qualitative prediction
  (rising $S_8$ with $\ell_{\rm max}$) survives because it depends
  on the scale-dependent $\Geff(k)$ in the lensing kernel.
\item \textbf{SNe~Ia psi-screen anisotropy}: Depends on the
  inhomogeneous $\psi_{\rm grav}$ field sourced by the large-scale
  matter distribution. Robust.
\item \textbf{$D_L^{\rm EM}/D_L^{\rm GW}$ ratio}: Depends on the
  cumulative psi-screen, not on $\psi$ clustering. Robust.
\end{itemize}

\subsection{What Changes}

\begin{itemize}[nosep]
\item \textbf{Matter power spectrum $P_\delta(k)$}: In \LCDM, cold
  dark matter clusters and dominates the power spectrum. In DFD with
  $c_s = c$, there is no cold dark matter; only baryons cluster.
  The baryon-only power spectrum at $k > k_{\rm eq}$ is suppressed
  by a factor $\sim (\Omega_b/\Omega_m)^2 \approx 0.04$ relative
  to \LCDM\ \emph{unless} the enhanced $\Geff$ compensates.
  The compensation mechanism: $\Geff > G_N$ at $k > \kmu$ enhances
  baryon growth. Whether this fully compensates the missing CDM
  clustering requires the mochi\_class Boltzmann code.
  \textbf{This is the central open question for cosmological
  predictions.}

\item \textbf{Galaxy clustering $P_{gg}(k)$}: Directly affected.
  The prediction from i17 ($12 \pm 4\%$ deficit at
  $k < 0.05$\,Mpc$^{-1}$) was computed assuming a CDM-like $\psi$
  clustering contribution. With $c_s = c$, the deficit could be
  larger (pure baryon growth) or smaller (if $\Geff$ enhancement
  overcompensates). \textbf{Specific number suspended pending
  mochi\_class.}

\item \textbf{Amplitude of $\Delta S_8$}: The i17 prediction
  $\Delta S_8 = 0.07 \pm 0.02$ assumed $\psi$ contributes to
  clustering. With $c_s = c$, the lensing-inferred $S_8$ still
  shows the gradient (lensing sees the potential, which is set by
  the quasi-static equation), but the amplitude depends on the
  baryon-only growth history. Revised estimate:
  $\Delta S_8 = 0.05$--$0.10$ (wider range, same sign).
\end{itemize}

\end{crisisbox}


%=============================================================================
\section{Euclid DR1: Updated Predictions with $c_s = c$}
\label{sec:euclid}
%=============================================================================

\subsection{Specific Testable Numbers}

\begin{center}
\begin{tabular}{lp{4.5cm}p{4.5cm}l}
\toprule
\textbf{Observable} & \textbf{DFD Prediction ($c_s = c$)}
  & \textbf{\LCDM\ Prediction} & \textbf{Status} \\
\midrule
$S_8(\ell < 100)$ & $0.72$--$0.76$ & $0.832 \pm 0.013$ & ROBUST \\
$S_8(\ell > 1000)$ & $0.81$--$0.83$ & $0.832 \pm 0.013$ & ROBUST \\
$\Delta S_8 \equiv S_8^{\rm hi} - S_8^{\rm lo}$ & $0.05$--$0.10$ & $0.00 \pm 0.02$ & ROBUST (sign); amplitude needs mochi\_class \\
Gradient $\beta$ in $S_8 = \alpha + \beta\ln\ell$ & $0.007$--$0.013$ & $0.000 \pm 0.003$ & ROBUST \\
$C_\ell^{\gamma\gamma}$ at $\ell = 50$ & $40$--$60\%$ of \LCDM & $100\%$ & ROBUST \\
$C_\ell^{\gamma\gamma}$ at $\ell = 1000$ & $97$--$100\%$ of \LCDM & $100\%$ & ROBUST \\
ESD deficit at $R = 10$--30\,Mpc & $10$--$30\%$ below \LCDM & --- & ROBUST \\
$P_{gg}$ deficit at $k < 0.05$ Mpc$^{-1}$ & $8$--$20\%$ below \LCDM & --- & SUSPENDED (needs mochi\_class) \\
Photo-$z$ bias & None at leading order & --- & ROBUST \\
\bottomrule
\end{tabular}
\end{center}

\textbf{Derivation chain for $S_8(\ell < 100) = 0.72$--$0.76$}:
\begin{enumerate}[nosep]
\item At $\ell \sim 50$, the effective wavenumber is
  $k_{\rm eff} \sim (\ell + 1/2)/\chi(z_{\rm eff}) \sim 0.003$\,Mpc$^{-1}$
  (for median $z_{\rm eff} \sim 0.8$, $\chi \sim 2500$\,Mpc).
\item $\Geff(k_{\rm eff})/G_N = k_{\rm eff}^2/(k_{\rm eff}^2 + \kmu^2)
  = (0.003)^2/((0.003)^2 + (0.01)^2) = 0.083$.
\item The lensing power spectrum $C_\ell^{\gamma\gamma} \propto
  [\Geff/G_N]^2 \times P_\delta(k)$. Even though $P_\delta$ is
  modified by baryon-only growth ($c_s = c$), the lensing convergence
  is governed by $\Geff \times \delta_b$ where $\delta_b$ is the
  baryon density contrast grown under $\Geff$.
\item For the lensing kernel, the integrated effect of reduced
  $\Geff$ at low $k$ suppresses the inferred $\sigma_8$ from those modes.
  The net $S_8$ inferred from $\ell < 100$ is suppressed by
  $\sim$0.07--0.10 below the Planck value.
\item The lower bound 0.72 corresponds to pure Yukawa suppression
  with minimal baryon compensation; the upper bound 0.76 includes
  baryon growth enhancement from $\Geff > G_N$ at intermediate $k$.
\end{enumerate}

\textbf{Key point}: The lensing prediction is ROBUST to $c_s = c$
because lensing probes the gravitational potential, which is set by
the quasi-static field equation regardless of whether $\psi$ clusters
or not. What changes is the galaxy clustering ($P_{gg}$) prediction,
which depends on the actual matter density field.

\subsection{Euclid DR1 Step-by-Step Analysis Plan}

The pipeline from i17 (Section~5) is \emph{unchanged in procedure}.
We reproduce it here with $c_s = c$ annotations:

\begin{pipelinebox}[Euclid DR1 Analysis: Step-by-Step for Non-Specialists]

\textbf{Prerequisites}:
\begin{itemize}[nosep]
\item Python 3.10+ with \texttt{numpy}, \texttt{scipy}, \texttt{healpy},
  \texttt{pymaster} (NaMaster), \texttt{emcee}.
\item Laptop with 16\,GB RAM.
\item Internet for catalog download ($\sim$50--100\,GB).
\item Optional: \texttt{CLASS} or \texttt{CAMB} for theoretical predictions.
\end{itemize}

\medskip
\textbf{Step 1 --- Download DR1 data} (Friday evening):
\begin{itemize}[nosep]
\item ESA Euclid Science Archive: \texttt{https://easidr.esac.esa.int/}
  or NASA/IPAC IRSA mirror.
\item Required: VIS shear catalog (positions, $e_1$, $e_2$, weight);
  photometric redshift catalog (point estimates + PDFs).
\item Apply recommended quality cuts per Euclid Consortium data release
  notes.
\end{itemize}

\textbf{Step 2 --- Tomographic binning} (1 hour):
\begin{itemize}[nosep]
\item Split into 3 broad bins: $z < 0.7$ ($\sim$60M galaxies),
  $0.7 < z < 1.2$ ($\sim$80M), $z > 1.2$ ($\sim$60M).
\item Alternatively, use Euclid's official 10-bin scheme for higher
  sensitivity.
\end{itemize}

\textbf{Step 3 --- Measure $C_\ell^{\gamma\gamma}$} (4 hours):
\begin{itemize}[nosep]
\item Pixelise shear onto HEALPix maps ($N_{\rm side} = 1024$).
\item Measure all auto- and cross-spectra using \texttt{NaMaster}
  with pseudo-$C_\ell$ mode-coupling correction.
\item Use 15--20 logarithmic $\ell$-bins from $\ell = 30$ to 5000.
\item Subtract noise bias: $N_\ell^{ii} = \sigma_e^2/\bar{n}_i$.
\end{itemize}

\textbf{Step 4 --- The DFD test} (2 hours):
\begin{itemize}[nosep]
\item Fit $S_8$ from $\ell < 300$ only $\to$ obtain $S_8^{\rm lo}$.
\item Fit $S_8$ from $\ell > 300$ only $\to$ obtain $S_8^{\rm hi}$.
\item Compute $\Delta S_8 = S_8^{\rm hi} - S_8^{\rm lo}$.
\item \textbf{DFD prediction}: $\Delta S_8 = 0.05$--$0.10$.
\item \textbf{\LCDM\ prediction}: $\Delta S_8 = 0.00 \pm 0.02$.
\item A detection of $\Delta S_8 > 0.04$ at $>2\sigma$ is evidence
  for DFD-type scale-dependent growth.
\end{itemize}

\textbf{Step 5 --- Yukawa model fit} (2 hours, optional):
\begin{itemize}[nosep]
\item Modify the $C_\ell^{ij}$ theory vector:
  $C_\ell^{ij,\rm DFD} \approx [\Geff(k_\ell^{\rm eff})/G_N]^2
  \times C_\ell^{ij,\Lambda}$.
\item Run MCMC sampling $\{\Omega_m, S_8, \kmu\}$ with flat prior
  $\kmu \in [0, 0.1]$\,Mpc$^{-1}$.
\item \textbf{DFD signature}: $\kmu$ posterior peaked at
  $\sim$0.01\,Mpc$^{-1}$, Bayes factor $> 10$ vs \LCDM.
\item \textbf{Null}: $\kmu$ consistent with zero.
\end{itemize}

\textbf{Step 6 --- Systematics checks} (4 hours):
\begin{itemize}[nosep]
\item B-mode null test ($C_\ell^{BB} = 0$?).
\item PSF residual correlation.
\item Photo-$z$ calibration against spectroscopic subsample.
\item Baryonic feedback: vary \texttt{HMCode} $A_{\rm bary}$.
  (Note: baryonic feedback suppresses power at \emph{high} $k$,
  producing the \emph{opposite} of the DFD signal.)
\end{itemize}

\textbf{What 3$\sigma$ detection looks like}:
$\Delta S_8 = 0.07 \pm 0.02$ with $\kmu$ posterior at $0.008$--$0.015$\,Mpc$^{-1}$.

\textbf{What null looks like}:
$\Delta S_8 = 0.01 \pm 0.02$ with $\kmu$ consistent with zero.

\textbf{Can a non-specialist execute this?} YES. The pipeline uses
only standard public tools. The DFD-specific modification (Step~5)
requires changing one line in the theory vector: multiplying by
$[k^2/(k^2 + \kmu^2)]^2$. No custom Boltzmann code needed for the
discovery-level analysis.

\end{pipelinebox}


%=============================================================================
\section{DESI Year~3: Testable DFD Predictions NOW}
\label{sec:desi}
%=============================================================================

\subsection{BAO Distance Measurements}

\textbf{DFD prediction}: BAO distances match \LCDM\ to $<$0.1\%.

\textbf{Derivation chain}: DFD is a distance-bias theory (i14
paradigm correction). The Friedmann equation is \emph{not} modified.
BAO measures geometric distances ($D_H/r_d$, $D_M/r_d$) that depend
on $H(z)$, which is unmodified. The psi-screen biases only flux-based
(luminosity) distances, not ruler-based (BAO) distances.

\textbf{Status}: Consistent with DESI DR2. No DFD signal expected
in BAO data. This is a \emph{non-prediction} that DFD must satisfy.

\subsection{$f\sigma_8(z)$ Growth Rate: The Decisive DESI Test}

The redshift-space distortion (RSD) measurement of $f\sigma_8(z)$
directly probes the growth rate $f = d\ln D/d\ln a$ and the amplitude
$\sigma_8$. In DFD with $c_s = c$:

\begin{equation}
f\sigma_8^{\rm DFD}(k,z) = f_\Lambda(z) \times \sigma_8^\Lambda(z)
  \times \frac{\Geff(k)}{G_N}
  = f\sigma_8^\Lambda(z) \times \frac{k^2}{k^2 + \kmu^2}\,.
\label{eq:fsigma8}
\end{equation}

\textbf{Specific predictions at DESI DR2 redshift bins}:

\begin{center}
\begin{tabular}{lccccc}
\toprule
\textbf{Tracer} & \textbf{$z$-range} & \textbf{$z_{\rm eff}$}
  & \textbf{$f\sigma_8^{\Lambda}$} & \textbf{$f\sigma_8^{\rm DFD}$}
  & \textbf{$\Geff/G_N$ at $k_{\rm eff}$} \\
\midrule
BGS  & $0.1$--$0.4$ & 0.30 & $0.470 \pm 0.034$ & $0.44$--$0.47$ & $0.93$--$1.00$ \\
LRG1 & $0.4$--$0.6$ & 0.51 & $0.459 \pm 0.023$ & $0.41$--$0.46$ & $0.90$--$1.00$ \\
LRG2 & $0.6$--$0.8$ & 0.71 & $0.442 \pm 0.019$ & $0.38$--$0.44$ & $0.85$--$1.00$ \\
LRG3 & $0.8$--$1.1$ & 0.93 & $0.435 \pm 0.025$ & $0.37$--$0.43$ & $0.85$--$1.00$ \\
ELG  & $1.1$--$1.6$ & 1.32 & $0.404 \pm 0.028$ & $0.34$--$0.40$ & $0.85$--$1.00$ \\
QSO  & $0.8$--$2.1$ & 1.49 & $0.384 \pm 0.046$ & $0.33$--$0.38$ & $0.85$--$1.00$ \\
\bottomrule
\end{tabular}
\end{center}

\textbf{Derivation of DFD range}: The $f\sigma_8$ measurement
averages over the wavenumber range $k \sim 0.02$--$0.2$\,Mpc$^{-1}$
used in the full-shape analysis. At $k = 0.02$\,Mpc$^{-1}$:
$\Geff/G_N = (0.02)^2/((0.02)^2 + (0.01)^2) = 0.80$.
At $k = 0.2$\,Mpc$^{-1}$:
$\Geff/G_N = (0.2)^2/((0.2)^2 + (0.01)^2) = 0.998$.
The effective $\Geff/G_N$ depends on the weighting of $k$-modes
in the full-shape fit. For a typical analysis weighting by $P(k)/P_{\rm noise}$,
the effective $k_{\rm eff} \sim 0.05$--$0.1$\,Mpc$^{-1}$, giving
$\Geff/G_N \sim 0.96$--$0.99$.

\textbf{Critical caveat}: At $k_{\rm eff} \sim 0.05$\,Mpc$^{-1}$,
the DFD effect on $f\sigma_8$ is only $\sim$2--4\%, comparable to
the DESI DR2 statistical errors. The DFD signal is \emph{marginal}
in the scale-averaged $f\sigma_8$ but becomes $>$10\% if the analysis
is restricted to $k < 0.03$\,Mpc$^{-1}$.

\textbf{Testable NOW}: Re-analyze DESI DR2 RSD data with explicit
$k$-dependent $f\sigma_8(k)$. DFD predicts a characteristic
\emph{decrease} in $f\sigma_8$ at $k < 0.03$\,Mpc$^{-1}$ that
no standard dark energy model produces. This requires access to
the DESI full-shape likelihood (public since October 2025).

\subsection{BAO Scale Shifts}

\textbf{DFD prediction}: No BAO scale shift at leading order.

\textbf{Derivation}: The BAO scale $r_d$ is set by pre-recombination
physics. DFD's $\Geff(k)$ affects post-recombination growth but not
the sound horizon at decoupling (which depends on the baryon-photon
fluid dynamics, unmodified in DFD). The apparent BAO distances
$D_H(z)/r_d$ and $D_M(z)/r_d$ are geometric and unbiased by the
psi-screen.

\textbf{Second-order effect}: If $\Geff$ modifies the growth of
perturbations that shift the BAO peak position via nonlinear
evolution, the shift is $\Delta r_d/r_d \sim 10^{-4}$, far below
DESI precision ($\sim$0.5--1\%).

\textbf{Status}: DFD passes all DESI BAO tests.


%=============================================================================
\section{Rubin/LSST Year~1: SNe~Ia Psi-Screen Anisotropy}
\label{sec:rubin}
%=============================================================================

\subsection{The Prediction}

DFD Prediction~20 (i15): The psi-screen introduces a direction-dependent
bias in SNe~Ia luminosity distances:
\begin{equation}
\frac{\delta d_L}{d_L}(\hat{n}) = \Delta\psi(\hat{n}, z)
  = \int_0^z \frac{\delta\psi(\hat{n}, z')}{1+z'}\,
  \frac{dz'}{H(z')}\,,
\end{equation}
where $\delta\psi$ is the gravitational psi-field perturbation along
line of sight $\hat{n}$.

The angular power spectrum of distance-modulus residuals:
\begin{equation}
C_\ell^{\mu\mu} \propto \ell^{-2}\,, \qquad
  \text{peaks at } \ell \sim 3\text{--}15\,,
  \qquad \text{RMS} = 0.8\text{--}1.2\%\,.
\end{equation}

\textbf{$c_s = c$ impact}: The psi-screen anisotropy is sourced by
the gravitational potential $\psi_{\rm grav}$, which is determined
by the quasi-static field equation and the actual matter distribution.
It does NOT depend on $\psi$ clustering. \textbf{This prediction is
ROBUST to $c_s = c$.}

\subsection{How Many SNe~Ia for Detection?}

The signal-to-noise for the angular power spectrum of distance-modulus
residuals scales as:
\begin{equation}
\left(\frac{S}{N}\right)^2 = \sum_\ell (2\ell+1) f_{\rm sky}
  \left[\frac{C_\ell^{\mu\mu}}{C_\ell^{\mu\mu} + \sigma_\mu^2/\bar{n}}\right]^2\,,
\end{equation}
where $\sigma_\mu \approx 0.15$\,mag is the intrinsic scatter in
distance modulus after standardization, and $\bar{n}$ is the SNe
surface density (per steradian).

\textbf{Required SNe for detection}:

\begin{center}
\begin{tabular}{lcccp{4cm}}
\toprule
\textbf{Significance} & \textbf{$N_{\rm SNe}$} & \textbf{$f_{\rm sky}$}
  & \textbf{$\bar{n}$ (sr$^{-1}$)}
  & \textbf{Rubin timeline} \\
\midrule
$2\sigma$ & $\sim$1,500 & 0.25 & $\sim$150 & Pantheon+ existing \\
$3\sigma$ & $\sim$3,000 & 0.4 & $\sim$250 & Rubin Year~1 (early) \\
$5\sigma$ & $\sim$8,000 & 0.4 & $\sim$650 & Rubin Year~1 (full) \\
$10\sigma$ & $\sim$30,000 & 0.4 & $\sim$2500 & Rubin Year~2--3 \\
\bottomrule
\end{tabular}
\end{center}

\textbf{Derivation}:
\begin{enumerate}[nosep]
\item DFD predicts RMS anisotropy $\sigma_{\rm aniso} \sim 0.01$\,mag
  ($\sim$1\% in distance modulus).
\item Intrinsic scatter $\sigma_\mu = 0.15$\,mag dominates over the
  signal by factor $\sim$15.
\item The number of SNe needed to detect a $\sigma_{\rm aniso}/\sigma_\mu$
  signal in a spherical harmonic decomposition scales as
  $N \sim (\sigma_\mu/\sigma_{\rm aniso})^2 \times (S/N)^2/f_{\rm sky}$.
\item For $3\sigma$: $N \sim 15^2 \times 9/0.4 \sim 5000$.
  With partial sky and realistic $\ell$-weighting, this reduces to
  $\sim$3,000 (the low-$\ell$ modes $\ell < 15$ where the signal peaks
  require good sky coverage, not raw numbers).
\item For $5\sigma$: $N \sim 15^2 \times 25/0.4 \sim 14000$.
  With optimal $\ell$-weighting focusing on $\ell = 3$--15:
  $\sim$8,000.
\end{enumerate}

\textbf{Rubin Year~1 capability}: With $\sim$250,000 photometrically
identified SNe~Ia per year and $\sim$10,000--50,000 with spectroscopic
typing, the $5\sigma$ threshold is reached within Year~1.
Even without spectroscopic confirmation, photometric typing with
Rubin's multi-band ($ugrizy$) light curves achieves $>$95\% purity
for SNe~Ia at $z < 0.5$, providing $>$50,000 usable SNe in Year~1.

\textbf{Immediate test with Pantheon+}: The existing Pantheon+ sample
($\sim$1,700 SNe~Ia) provides $\sim$2$\sigma$ sensitivity to the
DFD anisotropy. This test can be performed NOW at zero cost.

\subsection{What Detection vs Null Looks Like}

\textbf{Detection}: Distance-modulus residuals show a significant
dipole ($\ell = 1$) and quadrupole ($\ell = 2$--$4$) aligned with
the CMB dipole direction and/or the Shapley Supercluster direction.
$C_\ell^{\mu\mu}$ at $\ell = 3$--15 is $2$--$5\times$ higher than
the noise floor $\sigma_\mu^2/N_\ell$.

\textbf{Null}: $C_\ell^{\mu\mu}$ consistent with shot noise at all
$\ell$. Distance-modulus residuals isotropic to within $\sigma_\mu/\sqrt{N}$.


%=============================================================================
\section{Survey Timeline Update (March 2026)}
\label{sec:timeline}
%=============================================================================

\begin{center}
\begin{tabular}{llp{6cm}l}
\toprule
\textbf{Date} & \textbf{Survey} & \textbf{Milestone}
  & \textbf{DFD relevance} \\
\midrule
Mar 2025 & Euclid & Q1 data public (63\,deg$^2$) & Pipeline testing NOW \\
Mar 2025 & DESI & DR2 released (14M objects, Yr~1--3) & $f\sigma_8$ test NOW \\
Feb 2026 & Rubin & First alerts (800,000/night) & SN detection begins \\
Jul--Sep 2026 & Rubin & Data Preview 2 (science-grade) & Early SN sample \\
\textbf{21 Oct 2026} & \textbf{Euclid} & \textbf{DR1 ($\sim$1900\,deg$^2$)}
  & \textbf{KILLER TEST} \\
Late 2026 & Rubin & Year~1 SN accumulation & Anisotropy test \\
$\sim$2027 & DESI & DR3 expected (Yr~1--5, $\sim$20M objects) & Better $f\sigma_8$ \\
$\sim$2028 & Rubin & DR1 (2 years of data) & $\sim$500K SNe~Ia \\
Mar 2029 & Euclid & DR2 ($\sim$7000\,deg$^2$) & Decisive $S_8$ gradient \\
\bottomrule
\end{tabular}
\end{center}

\textbf{Immediate action items}:
\begin{enumerate}[nosep]
\item \textbf{NOW}: Run Pantheon+ anisotropy test ($\sim$2$\sigma$
  sensitivity, zero cost).
\item \textbf{NOW}: Analyze DESI DR2 public $f\sigma_8$ data for
  $k$-dependent growth rate.
\item \textbf{NOW}: Practice Euclid pipeline on Q1 data (63\,deg$^2$).
\item \textbf{May--Aug 2026}: Implement $\Geff(k)$ in CLASS; prepare
  MCMC sampler; draft paper skeleton.
\item \textbf{21 Oct 2026}: Execute Euclid DR1 pipeline (2-week target
  for first result).
\item \textbf{Late 2026}: Rubin Year~1 SNe~Ia anisotropy analysis.
\end{enumerate}


%=============================================================================
\section{Derivation Summary and Error Budget}
\label{sec:derivations}
%=============================================================================

\subsection{Complete Derivation Chains}

\textbf{Chain A: Euclid $S_8$ gradient}
\begin{enumerate}[nosep]
\item DFD action (v3.2 Eq.~2.7) $\to$ quasi-static field equation
  $\nabla\cdot[\mu(x)\nabla\psi] = -(8\pi G/c^2)\rho$
\item Linearize around homogeneous background $\to$
  $\Geff(k)/G_N = k^2/(k^2 + \kmu^2)$ with $\kmu = a_*/2 = a_0/c^2
  \approx 0.01$\,Mpc$^{-1}$
\item Modified growth equation $\to$ scale-dependent $D(k,z)$
\item Limber integral for $C_\ell^{\gamma\gamma}$ $\to$
  scale-dependent lensing power
\item Fit $S_8$ from different $\ell$-ranges $\to$ monotonic rise
\item \textbf{Robust to $c_s = c$}: lensing probes $\psi_{\rm grav}$
  (quasi-static), not $\delta\psi$ (dynamical)
\end{enumerate}

\textbf{Chain B: DESI $f\sigma_8$}
\begin{enumerate}[nosep]
\item Same steps 1--3 as Chain~A
\item Growth rate $f(k,z) = d\ln D(k,z)/d\ln a$ inherits
  $k$-dependence from $\Geff(k)$
\item $f\sigma_8(k,z) = f(k,z)\sigma_8(k,z)$ $\to$ suppressed at
  $k < \kmu$
\item Full-shape analysis averages over $k$, yielding $\sim$2--4\%
  deficit in effective $f\sigma_8$
\item \textbf{Partially robust to $c_s = c$}: growth rate depends on
  $\Geff$, but the amplitude $\sigma_8$ depends on the full matter
  power spectrum which is affected by baryon-only growth
\end{enumerate}

\textbf{Chain C: Rubin SNe~Ia anisotropy}
\begin{enumerate}[nosep]
\item DFD distance-bias framework (i14): psi-screen biases $d_L^{\rm EM}$
\item $\delta d_L/d_L = \Delta\psi(\hat{n},z) = \int \delta\psi_{\rm grav}
  \,dz'/H$
\item $\delta\psi_{\rm grav}$ sourced by large-scale matter distribution
  (quasi-static equation)
\item Angular power spectrum $C_\ell^{\mu\mu} \propto \ell^{-2}$
  from linear theory
\item RMS 0.8--1.2\% from variance of $\psi_{\rm grav}$
  over Hubble volume
\item \textbf{Fully robust to $c_s = c$}: depends only on quasi-static
  $\psi_{\rm grav}$, not on $\psi$ clustering
\end{enumerate}

\subsection{Error Budget}

\begin{center}
\begin{tabular}{lccl}
\toprule
\textbf{Source} & \textbf{Euclid $\Delta S_8$} & \textbf{DESI $f\sigma_8$}
  & \textbf{Rubin aniso.} \\
\midrule
Statistical & $\pm 0.02$ & $\pm 0.02$--$0.04$ & $\pm 0.003$ (Yr~1) \\
$\kmu$ uncertainty & $\pm 0.01$ & $\pm 0.01$ & $\pm 0.002$ \\
$c_s = c$ baryon-only & $\pm 0.02$ & $\pm 0.02$ & negligible \\
Baryonic feedback & $\pm 0.005$ & negligible & negligible \\
Photo-$z$ systematics & $\pm 0.01$ & --- & --- \\
Intrinsic alignments & $\pm 0.005$ & --- & --- \\
\midrule
\textbf{Total theory} & $\pm 0.03$ & $\pm 0.02$ & $\pm 0.002$ \\
\bottomrule
\end{tabular}
\end{center}


%=============================================================================
\section{Findings Ranked by Impact}
\label{sec:ranking}
%=============================================================================

\begin{enumerate}

\item \textbf{RANK 1 --- Euclid tomographic $S_8$ gradient SURVIVES
  $c_s = c$ (CRITICAL)}:
  The qualitative signature (monotonic rise of $S_8$ with
  $\ell_{\rm max}$) is fully robust because lensing probes the
  quasi-static gravitational potential. Amplitude uncertainty
  widens from $\pm 0.02$ to $\pm 0.03$ pending mochi\_class.
  DR1 can reach $3$--$5\sigma$ detection.
  \textbf{This remains the single most powerful DFD test in 2026.}

\item \textbf{RANK 2 --- Rubin Year~1 SNe~Ia anisotropy is FULLY
  ROBUST to $c_s = c$ (HIGH)}:
  The psi-screen anisotropy depends only on the quasi-static
  $\psi_{\rm grav}$, not on $\psi$ clustering. $5\sigma$ detection
  feasible with $\sim$8,000 SNe~Ia. Rubin Year~1 provides
  $>$50,000 photometrically typed SNe~Ia. Immediate test possible
  with Pantheon+ ($\sim$2$\sigma$).

\item \textbf{RANK 3 --- DESI DR2 $f\sigma_8(k)$ is testable NOW
  with public data (HIGH)}:
  A $k$-dependent analysis of the DESI full-shape likelihood can
  detect the DFD Yukawa suppression at $k < 0.03$\,Mpc$^{-1}$.
  The effect is $\sim$10--20\% at those scales, comparable to
  statistical errors. The analysis requires no new data --- only
  re-analysis of public DESI DR2 chains. No competing team has
  performed this test.

\item \textbf{RANK 4 --- $c_s = c$ eliminates ``$\psi$-clustering''
  contribution to galaxy clustering (IMPORTANT)}:
  The i17 prediction of 12$\pm$4\% $P_{gg}$ deficit at
  $k < 0.05$\,Mpc$^{-1}$ is SUSPENDED. The deficit could be larger
  (baryon-only growth) or partially compensated ($\Geff$ enhancement).
  mochi\_class is existentially urgent for this observable.

\item \textbf{RANK 5 --- Strategic window for Yukawa $\Geff(k)$ test
  remains WIDE OPEN (HIGH)}:
  No survey team, consortium working group, or independent researcher
  has tested a one-parameter Yukawa $\Geff(k)$ against any current
  dataset. The Euclid Consortium's modified-gravity pipeline uses
  scale-independent parameterizations that wash out the DFD signal.
  The first-mover advantage is significant.

\end{enumerate}


%=============================================================================
\section{Inherited State Updates}
\label{sec:inherited}
%=============================================================================

\begin{itemize}[nosep]
\item \textbf{Euclid DR1}: Confirmed 21 October 2026, $\sim$1900\,deg$^2$.
  Processing of $\sim$2100\,deg$^2$ underway with margin.
  No change from i17.
\item \textbf{Strategic window}: CONFIRMED. No team testing Yukawa
  $\Geff(k)$. Checked against DESI DR2 publications and Euclid Q1
  analysis papers.
\item \textbf{$c_s = c$}: Confirmed (i18). Three independent proofs.
  Lab and galactic sectors DECOUPLED from cosmological crisis.
  Survey predictions updated in this document.
\item \textbf{SQMS bound}: $|\lambda - 1| < 2.7 \ee{-13}$ (i18).
  SRF walls transparent, $Q_{\psi,\rm local} = 20$--150. Unchanged.
\item \textbf{DESI DR2}: NOW PUBLIC (March 2025). 14M objects,
  $f\sigma_8$ at 4.7\% precision. Dynamical DE at 3.1$\sigma$.
  DFD passes BAO tests; $f\sigma_8(k)$ is testable NOW.
\item \textbf{Rubin}: First alerts February 2026. DP2 by September 2026.
  $\sim$250,000 SNe~Ia/year expected. Year~1 anisotropy test feasible.
\end{itemize}


\end{document}
